\begin{abstract}
Los robots aspiradora, los vehículos autónomos y los equipos de topografía
construyen mapas a partir de un problema común: un sensor mide distancias
respecto a sí mismo y esas medidas deben llevarse a un sistema de referencia
compartido. Este trabajo presenta un escáner LIDAR 2D de bajo costo, construido
con un microcontrolador ESP32, un sensor láser de tiempo de vuelo VL53L0X y un
servomotor SG90, que resuelve ese problema con álgebra lineal. Cada lectura
$(r,\theta)$ se convierte en un punto del plano mediante la matriz de rotación
$R(\theta)$, y los escaneos tomados desde dos posiciones se fusionan en un solo
mapa con una matriz homogénea $3\times3$ que combina rotación y traslación en
un único producto matricial. Una interfaz web dibuja el mapa en tiempo real y
muestra la multiplicación de matrices de cada lectura. En un entorno simulado,
el mapa fusionado reunió \cifPuntosTotal{} puntos y el ruido del sensor produjo
un error medio de \cifRuidoMedio{} mm. Un error de $5^\circ$ en el ángulo de la
segunda posición elevó el error medio a \cifErrorCinco{} mm, en concordancia
exacta con la expresión teórica $2|\sin(\delta/2)|\,\|\vect{v}\|$, lo que
muestra que la orientación del aparato pesa más que su desplazamiento en la
calidad del mapa. Con el prototipo real se capturaron dos barridos completos, y la
distancia medida en un mismo ángulo varió en promedio \realRepMedia{} mm
entre ellos. Se verificó numéricamente que las matrices implementadas son
ortogonales, conservan la norma y tienen inversa exacta.
\end{abstract}

\keywords{álgebra lineal, matriz de rotación, coordenadas homogéneas,
cambio de base, LIDAR, mapeo de entornos}
